% ---------------------------------------------------------------------
% Das Dokument kompiliert mit pdflatex und ist auf Basis 
% von Koma-Script entstanden. 
%
% Autor des Templates (für Anmerkungen): 
% Michael von Riegen, riegen@informatik.uni-hamburg.de
%
% Einzelne Code-Teile für das Titelblatt sind aus dem Template 
% von Benjamin Kirchheim entnommen.
%
% 25.05.09, Frank Langanke: Vorlage auf aktuelle KOMA-Version aktualisiert
% 26.05.09, Michael von Riegen: Anmerkung --> aktuelles Koma-Script ist nötig!
% 17.10.2016 neues Uni logo
% ---------------------------------------------------------------------
\documentclass[11pt,DIV=15,BCOR=20mm,bibliography=totoc]{scrbook}

% Import von Paketen und Optionen die das gesamte Dokument betreffen
% sind in myPreamble.sty ausgelagert.
\usepackage{myPreamble}
   
% Arbeitet man nur an einem Kapitel, wird durch folgenden Befehl nur dieses eingebunden.
% Spart manuelles auskommentieren von vielen include-Befehlen;
% hat keine Auswirkung auf input-Befehle
% \includeonly{kapitel1}
   
\begin{document}

% TITELSEITE
\begin{titlepage}

	% Fehler "destination with the same identifier" unterdrücken...
  \setcounter{page}{-1}

	% Titelseite
	\begin{figure}[h]
		\begin{minipage}[b]{62mm}
			\includegraphics[width=62mm]{images/unilogo}
		\end{minipage}
		\hspace{4cm}
		\begin{minipage}[b]{59mm}
			\includegraphics[width=59mm]{images/minlogo}
		\end{minipage}
	\end{figure}

	\vfill
	
	\begin{center}
		
		\noindent { \large{}
			Hausarbeit \\
		}
		\vspace{14mm}
		% Titel
		\noindent \huge{\
		  Robotik in der Altenpflege: Ethische Überlegungen zur sozialen Interaktion.
		}
		\vspace{60mm}	
	\end{center}
	
	\vfill
	
	\noindent \textbf{Melanie Budde} \\
	\noindent \rule{\textwidth}{0.4mm} 
	\noindent{\textrm{2budde@informatik.uni-hamburg.de}} \\
	\noindent{\textrm{Studiengang Informatik}} \\
	\noindent{\textrm{Matr.-Nr. 6494505}} \\

	%\begin{tabbing}
	%\hspace{8em} \=  \kill

    %\noindent{\textrm{Erstgutachter: Dr. Julian Kunkel}} \\
	%\noindent{\textrm{Zweitgutachter: Prof. Dr. Thomas Ludwig}} \\
	%\Erstgutachter: \> Dr. Julian Kunkel \\
	%\Zweitgutachter: \> Prof. Dr. Thomas Ludwig \\
	 \\
%	Abgabe: 31.08.2019
%	\end{tabbing}
	
	% Rückseite der Titelseite mit Zitat
	\newpage 
	\thispagestyle{empty}
	\setcounter{page}{0}

	% wenn man Lust auf ein Zitat hat...
	% ... ansonsten auskommentieren
	%~\\ \vfill \noindent 
	%A distributed system is one where the failure of some \\
	%computer I've never heard of can keep me from getting my work done. \\
	%\textit{-- Leslie Lamport}
\end{titlepage}



% VERZEICHNISSE (Inhaltsverzeichnis, Abkürzungen)
% Vorspann einleiten --> Seitennummerierung römisch
\frontmatter

% Inhaltsverzeichnis
\tableofcontents
\cleardoublepage

% Hauptteil einleiten --> Seitennummerierung wieder arabisch
\mainmatter

\chapter{Einleitung}
\label{ch:einleitung}
Die Auseinandersetzung mit dem Thema Pflegerobotik hat in den letzten Jahren sowohl im öffentlichen Diskurs als auch in der wissenschaftlichen Literatur zugenommen. Dies liegt einerseits daran, dass das Gesundheitswesen mit Herausforderungen konfrontiert ist, für die eine innovative Herangehensweise am wahrscheinlichsten Lösungen verspricht. Zu den Herausforderungen zählen sowohl die demografische Entwicklung als auch der Mangel an Fachkräften. Andererseits werden von Herstellern und Universitäten vermehrt Projekte vorgestellt, welche einen Nutzen für die Praxis vorweisen können. Im folgenden wird der Einsatz sozial interagierender Roboter in der Altenpflege diskutiert. Dabei wird in \ref{sec:values} auf die Wechselwirkung zwischen Design und Werten eingegangen. In Abschnitt \ref{sec:täuschung} wird eine ethische Reflektion zum Thema der Täuschung im Bereich der Robotik dargelegt. Fragen der Privatheit, die diesbezüglich auftauchen werden in \ref{sec:priv} behandelt. Anschließend werden Bedenken bisheriger Ansätze zusammengefasst.



%Wie  kürzlich  in  einer  Forschungsarbeit  von  Bilyea  etal.  (2017)  ausgeführt,  haben  die  meisten   technologischen   Fortschritte   im   Bereich   der   Pflegeroboter   in   den   letzten   12Jahren  stattgefunden.  Es  ist  deshalb  wenig  überraschend,  dass  Forschung  in  diesem  Bereich  noch  in  den  Kinderschuhen  steckt.  %Frühe  Arbeiten  von  Scopelliti  etal.  (2005) zeigen,  dass  die  Akzeptanz  von  und  Einstellung  zu  Robotern,  insbesondere  bei  älteren  Menschen,  nicht  nur  von  den  praktischen  Vorteilen  abhängen,  sondern  aus  komplexen  Beziehungen  zwischen  den  kognitiven,  affektiven  und  emotionalen  Komponenten  der  Menschen  und  ihren  Vorstellungen  von  Robotern  entstehen.  Sie  fanden  heraus,  dass  ältere  Menschen  ein  höheres  Misstrauen  gegenüber  der  Technik  hatten  und  diese  auch  komplizierter  zu  bedienen  fanden.  In  ähnlichen  Studien  hat  sich  gezeigt,  dass  ältere  Menschen auch früher aufgaben, wenn sie Schwierigkeiten hatten (Giuliani etal. 2005). Andererseits  stellten  Stafford  etal.  (2014)  fest,  dass  Bewohner  von  Alters-  und  Pflege-heimen, die physische Einschränkungen hatten, einen Roboter eher benutzten als körper-lich gesunde Bewohner. Ähnlich haben Tinker und Lansley (2005) und Pain etal. (2007)dokumentiert, dass ältere Menschen besonders bereit waren, Technologie zu akzeptieren, wenn sie sich an einen spezifischen Bedarf richtete und ihnen somit mehr Unabhängig-keit brachte. Forschung, die als Kernpunkt die Interaktion zwischen Mensch und Robo-ter  hat,  war  bisher  meist  beobachtend,  unterlag  künstlichen  Bedingungen  und  wurde  nur  mit  Fotos  und  Videos  von  Robotern  durchgeführt,  statt  experimentell  und  in  rea-len  Umgebungen.  Daten  wurden  überwiegend  durch  Fragebogen  (Giuliani  etal.  2005; 
%48M. Früh und A. GasserScopelliti  etal. 2005),  Nachbearbeitungsinterviews  (Birks  etal.  2016)  oder  Fokus-gruppen  (Broadbent  etal.  2012;  Wu  und  Rigaud  2012)  gesammelt  anstelle  von  Inter-aktionen mit Robotern in der Zielu [3]%

%Bei  älteren  Menschen  ist  die  Akzeptanz  geringer  als  bei  jüngeren  Menschen.  Dennoch  zeigen  Umfragen,  dass  immer  mehr  Seniorinnen und Senioren dem Einsatz von Pflegerobotern positiv gegenüberstehen. Bei-spielsweise konnten sich bei einer repräsentativen Forsa-Umfrage, welche im April 2016 im Auftrag des BMBF durchgeführt wurde, 83 der Befragten grundsätzlich vorstellen, im Alter einen Serviceroboter in ihrem eigenen Zuhause zu nutzen, um dort länger woh-nen zu können (Forsa. Politik- und Sozialforschung GmbH 2016). [4]









\chapter{Ethische Erläuterungen zu Robotik in der Pflege}
\label{ch:2}
 Als Gründe für die Entwicklung von Robotern in der Pflege führt sie den demographischen Wandel in vielen Industrienationen, den Pflegenotstand, die Beförderung von Unabhängigkeit in der häuslichen Pflege durch Monitoring und Überwachung, sowie die Kostenreduzierung an \cite[56]{haker}.  Dabei beherbergt gerade die Entwicklung von Pflegerobotern noch viele ungeklärte ethische Hürden. Im folgenden wird sich lediglich mit sozial interagierenden Robotern in der Altenpflege  beschäftigt. Für Janowski unterscheiden sich sozial  interagierende  Roboter  von  reinen  Service-robotern  dadurch,  dass  sie  während  der  Ausführung  von  Diensten  Verhaltensweisen  emulieren,  die  an  der zwischenmenschlichen  Kommunikation  orientiert  sind. \cite[65]{vier}
Da dieser Bereich der Robotik noch vergleichsweise in den Kinderschuhen steckt, gibt es nur wenig Erfahrung mit dem Einsatz von Pflegerobotern und wo ein politisches Vakuum darüber herrscht, wie Computertechnologie eingesetzt werden sollte, herrscht häufig auch ein konzeptuelles \cite[1]{comEthics}.  Verschiedene Veröffentlichungen beschäftigen sich in den letzten Jahren mit der Frage, inwiefern der Einsatz von Robotik in der Pflege ethisch vertretbar ist. Haker \cite[56]{haker} nennt drei Bereiche, in denen neue Technologien hilfreich sein können und neuerdings erprobt werden: die Pflegeassistenz, das Monitoring sowie die soziale Förderung.



\section{Beziehung zwischen Design und Werten von Technologie}
\label{sec:values}
Es ist wichtig, frühzeitig ethische Probleme während des Designprozesses zu bedenken.
Dabei sagt Moore: "On my view, computer ethics is the analysis of the nature and social impact of computer technology and the corresponding formulation and justification of policies for the ethical use of such technology" \cite[1]{comEthics}. Eine Analyse der Wechselwirkungen zwischen Technologie und sozialen Konstrukten ist demnach unabdingbar, wenn man ein ethisch korrektes Computersystem designen möchte. Nissenbaum setzt sich hierbei mit der Frage auseinander, inwiefern \textit{values in design}, also die durch das jeweilige Design implizierten Werte, besser in technischen und naturwissenschaftlichen Bereichen reflektiert und integriert werden können \cite{nissenbaum}. Oftmals gibt es in diesen universitären Fakultäten wenig Überschneidungen mit den Sozial- und Geisteswissenschaften, welche eine längere Tradition der Methodologie aufweisen [ebd.: lxvi]. Sie schlägt vor, dass pflichtbewusste Designer*Innen ihre values in design auf drei Wissensebenen reflektieren. Zuerst werden diese hinsichtlich des Wissensstandes in Technik und Naturwissenschaft überprüft, folgend durch eine philosophische Reflexion, und zuletzt durch eine empirische Überprüfung im Anwendungsfeld [ebd.]. Durchgängig sollte dabei im Fokus stehen, dass sich alle drei Bereiche durchdringen und gegenseitig beeinflussen. Sie sollten also, während der Fokus sich in einem dieser Bereiche befindet, kontinuierlich mitgedacht werden. Während also im ersten Schritt sämtliche im jeweiligen technischen Feld vorhandenen Werte auf den Prüfstand gestellt werden um sich sinnvoll für ein Design zu entscheiden, wird im zweiten Schritt geschaut, wo Werte ihren Ursprung haben, welche Relevanzen und Bedeutungen ihnen innewohnen und wie sie normativ beeinflusst sind [ebd.: lxvii]. Werte können eine moralische, eine politische und eine soziale Konnotation haben [ebd.]. Werden Werte ins Design bewusst integriert, sollte im nächsten Schritt überprüft werden, ob ihre Einverleibung geglückt ist, also ob die Intentionen der*s Designer*In sich erfüllen [ebd.; lxviii].  Möglich ist, ein partizipatives Design zu erstellen, wo die Entwicklungsschritte einer Technologie oder einer Maschine jeweils rückgekoppelt sind an eine Überprüfung auf ihre Wirksamkeit im späteren Anwendungsbereich [ebd.]. Nissenbaum hat mit Kolleg*Innen hierfür eine eigene Methode entwickelt, in der hierfür vier Schritte durchlaufen werden \cite{nissenbaum}. Zunächst die Entdeckung der jeweils zugrunde liegenden Werte, folgend ihre Übersetzung in das jeweilige Feld, danach die Lösung zumeist auftretender Konflikte oder Inkontinenzen und zuletzt die Überprüfung einer erfolgreichen Einverleibung [ebd.: lxix]. Der Einsatz von Robotern in der Pflege allgemein, aber auch in Bezug auf soziale Interaktion, birgt viele Möglichkeiten, über eine reflektierte Herangehensweise bezüglich der values in design, die ethischen Vorraussetzungen positiv zu beeinflussen. Wie Sharkey und Sharkey und Kehl feststellen, bedarf es hier einer genauen Überprüfung der Grundvorraussetzungen und eine Zusammenarbeit mit den Pflegebedürftigen \cite{sharky}\cite{kehl}. Mithilfe empirischer Untersuchungen und  partizipativer Technikgestaltung wäre es möglich, wichtige Werte von Älteren herauszufinden und aufzugreifen und somit auf technischer Seite dafür zu sorgen, dass ethische Fragen in das Design eingearbeitet werden. 

\section{Kann Täuschung ethisch Vertretbar sein?}
\label{sec:täuschung}
\\Auch Sharkey und Sharkey sehen als wichtige Bereiche, in denen Roboter eingesetzt werden können, die Unterstützung der Älteren, sowie der sie pflegenden; das Monitoring von Verhalten und Gesundheitszustand; sowie das Gesellschaftleisten \cite[ 27]{sharky}. Menschen können durch Roboter emotional stimuliert werden und diese können dabei helfen “Gefühle wie Traurigkeit, Zorn und Einsamkeit zu regulieren” \cite[58]{haker}. Was die ethischen Konsequenzen angeht, stimmen die meisten Veröffentlichungen hinsichtlich der sich ergebenen Problematiken überein, erkennen jedoch auch den Nutzen an. An dieser Stelle wird sich eine Analyse der sozialen Unterstützung durch Roboter beschränkt. Haker stellt drei Modelle vor, in denen sie die unterschiedlichen Argumentationslinien bezüglich der Position von Robotern in sozialen Beziehungen im Pflegebereich darlegt \cite[59 ff]{haker}. Im Ausschlussmodell wird angenommen, dass zwischen Menschen und Robotern keine gleichwertige Beziehung existieren kann. Roboter werden als zu sozialer Interaktion nicht fähig gesehen, problematisch sei es, wenn Pflegebedürftige dieses durch kognitive Einschränkungen nicht mehr nachvollziehen können. “[…]So warnen Kritiker nicht nur vor der Infantilisierung der Klienten, die durch die BetreuerInnen oder Angehörigen erfolgen könnte, sondern zugleich besteht auch die Gefahr der Objektivierung, indem letztlich nämlich mit den Bedürfnissen älterer Menschen gespielt werde. Die absichtliche Simulation reziproker Interaktion durch Dritte ist von einer anderen Kategorie als die selbstbestimmte Entscheidung älterer Menschen, Roboter als soziales Spielzeug in ihr Leben einzubeziehen” [ebd.: 60]. Im Ergänzungsmodell werden Roboter vor allem für gewisse Arbeiten gedacht, die durch eine Ergänzung der Pflegeleistung diese unterstützen. Dabei wird der Einsatz von Robotern immer durch das Pflegepersonal betreut, “[…]ihr Einsatz ist klar durch den therapeutischen Rahmen gegeben, den weder die Klienten (bzw. Patienten) noch die Roboter durchbrechen können oder sollen” [ebd.: 61]. Haker wirft die Frage auf, ob in dieser Argumentationslinie das autonome Handlungsvermögen von Robotern nich unterschätzt wird \cite{haker}.\\ Im Kompensationsmodell geht es darum durch den Einsatz von Robotern die Abwesenheit anderer Interaktion zu kompensieren. Wenn soziale Interaktion mit Angehörigen oder Pflegepersonal zeitlich nicht möglich ist, sei es immer noch besser wenn Roboter den zu Pflegenden ein Wohlgefühl vermitteln und ihnen Gesellschaft leisten würden. Die Argumentationslinie des Kompensationsmodells bewertet das Steigern des Wohlbefindens als wichtiger, selbst wenn die zu Pflegenden bezüglich der Reziprozität der Beziehung getäuscht werden würden, falls sie kognitiv nicht mehr feststellen können, dass sie gerade mit einem Roboter interagieren [ebd.: 62].\\ Je nach dem, welcher Argumentationslinie gefolgt wird, wird die Annahme vertreten, dass soziale Beziehungen zu Robotern niemals gleichwertig sein können, was im schlimmsten Fall von Älteren oder Demenzkranken nicht gemerkt wird, womit sie also getäuscht werden. Oder der Einsatz von Robotern wird dennoch als positiv bewertet, da anderweitige soziale Interaktion bewusst ergänzt werden kann, oder im schlechteren Fall die Abwesenheit sozialer Interaktion durch diese andere Form der Interaktion kompensiert.\\Sparrow und Sparrow lehnen den Einsatz von Robotern nicht grundsätzlich ab, befinden ethische Bedenken jedoch als zentral. Sie weisen darauf hin, wie kompliziert die einzelnen Aufgaben im Bereich der Pflege sind, und bezweifeln, dass künstliche Intelligenz in naher Zukunft solch komplexe und verantwortungsvolle Aufgaben übernehmen werden kann \cite[145 ff]{sparrow}. Sparrow und Sparrow gehen davon aus, dass das zentrale menschliche Bedürfnis, geliebt zu sein und umsorgt zu werden, von Robotern nur befriedigt werden kann, wenn man der Täuschung unterliegt, dass diese gleichwertige Kreaturen mit Gefühlen sind \cite[154 ff]{sparrow}. Beim Empfinden dieser positiven Emotionen ist die Beziehung zwischen zwei Lebewesen zentral. Sie ist dadurch charakterisiert, dass auch ein Gegenüber Bedürfnisse und Erwartungen hat, und dass gegenseitige Forderungen die Beziehung formen, welche sowohl unvorhersehbar, als auch manchmal anmaßend und unerwartet sein können [ebd.: 149]. “For robots to be capable even of imitating these responses successfully they would need to possess physical bodies capable of the same level of expressiveness and individuality as human bodies. Moreover, entities which do not understand the facts about human experience and mortality that make tears appropriate will be unable to fulfil this caring role” \cite[154]{sparrow}. Die Täuschung, sich von einem Roboter umsorgt und geliebt zu fühlen, ist vor allem problematisch, da Illusionen wirkliche Bedürfnisse nicht dauerhaft befriedigen, wirkliche Interessen nicht dauerhaft bedienen und auch ein Wohlgefühl nicht dauerhaft illusionär, oder durch Selbsttäuschung, aufrecht erhalten werden kann [ebd.:155]. Menschen wollen nicht in illusionärer Annahme glauben, dass sie von illusionären technischen Maschinen geliebt und umsorgt werden, sondern sie haben das reale Bedürfnis von Menschen geliebt und umsorgt zu werden [ebd.]. “Insofar as robots can make people happier only when they are deceived about the robots’ real nature, robots do not offer real improvements to people’s well-being; in fact the use of robots can be properly said to harm them“ [ebd.].  Auch Haker stellt diesbezüglich fest: “Ohne diese Form der imaginativen Personalisierung des Roboters verlöre das Spiel mit ihm schnell seinen Reiz – oder umgekehrt: es wäre nicht mehr klar, was der Unterschied beispielsweise zu Stofftieren wäre, mit denen ältere Menschen genauso imaginativ soziale Interaktionen aufbauen könnten“ \cite[59]{haker}. Sparrow und Sparrow vermuten, dass positive emotionale Effekte, die aus der Interaktion mit Robotern hervorgehen, sich nicht über einen langen Zeitraum halten lassen, da eine authentische Beziehung nicht aufgebaut werden kann \cite[149]{sparrow}. Diese Effekte beruhen somit bloß auf der Faszination einer neuen, anregenden Erfahrung. Sharkey und Sharkey beantworten die Frage, ob ein Roboter für Ältere ein*e Gefährt*In darstellen kann, mit der Festellung, dass einsame Menschen von der Anwesenheit eines Roboters profitieren können, aber sicherlich um einiges mehr profitieren würden, wenn sie sich in einem freundlichen sozialen Umfeld befänden \cite[34]{sharky}. Sie schreiben jedoch auch, dass Roboter als Gefährten nicht einfach aufgrund der stattfindenden Täuschung als umethisch gewertet werden können. Ältere könnten durchaus bewusst entscheiden, dass sie die Interaktion mit Robotern genießen, auch wenn es sich um eine illusionäre Beziehung handelt [ebd.: 36]. Auch von einer Infantilisierung kann nicht in jedem Fall gesprochen werden, es gibt durchaus Situationen wo ein älterer Mensch sich entscheidet einfach die Möglichkeit zu nutzen und das Beste daraus zu machen.
Das Aussehen von Robotern ist eine weitere wichtige Frage, werden diese eher dem Aussehen von Lebewesen nachempfunden, oder wird ein abstraktes Äußeres bevorzugt. Je abstrakter die Erscheinung eines Roboters ausfällt, desto eher lässt sich vermuten, dass Pflegebedürftiger einer Täuschung weniger schnell unterliegen. Es wäre möglich den Roboter so zu kreieren, dass einer Vereinsamung von Pflegebedürftigen durch Warnmechanismen entgegengewirkt wird. Werden diese im Hinblick darauf entwickelt, um Pflegebedürftigen zu helfen, oder im Hinblick darauf, Kosten im Pflegebereich einzusparen und das Arbeitspensum des Pflegepersonals zu reduzieren \cite[30]{sharky}.

\section{Probleme mit der Privatsphäre}
\label{sec:priv}
In der Literatur angesprochene Themen sind dabei, wie eine Kameraüberwachung in das Aussehen und Auftreten von Robotern integriert werden kann oder unter welchen Voraussetzungen gewisse private Räumen wie Bad und Schlafzimmer von Robotern betreten werden dürfen. Denn erst durch die Schaffung eines privaten Raumes, kann es zu der Ausübung von Autonomie kommen \cite[178]{priv}. Dieser wird aber beschnitten, sobald z.B dauerhafte Kameraüberwachung installiert wird. Somit ist fraglich, ob z.B aufzeichnende Assistenzsysteme in Privathaushalten wirklich die Autonomie fördern. Desweiteren gehen Grimm und Krah darauf ein, dass vier Funktionen für die Gewährleistung von Privatheit erfüllt sein müssen: 1. \textit{personal autonomy}, 2. \textit{emotional release}, 3. \textit{self-evaluation}, 4. \textit{limmited and protectet communication} [ebd.: 178]. Die Einführung eines Roboters der immer zu erreichbar ist, fördert zwar den ersten Punkt, verletzt jedoch den vierten.
Darüber hinaus werden bei jeder Interaktion Daten erzeugt, was eine Menge Fragen mit sich bringt: Wem gehören die Daten? Sollten sie gespeichert werden und wenn ist zu klären wo, Dürfen die Daten weiter verarbeitet werden? Dazu sagt Sullins: "In one sense this could be immensely powerful data that could lead to much healthier lifestyle choices. But it could also be a serious breach in privacy if the information got into the wrong hands which would be easily accomplished since third parties have access to information collected on smartphones and online applications." \cite{sullins}. 

\section{Bedenken bei der Einführung von Robotern}
Sharkey und Sharkey untersuchen den Einsatz von Robotern in der Pflege vor allem unter dem Kriterium des Wohlergehens älterer pflegebedürftiger Menschen \cite[S. 28]{sharky}. In Bezug auf die unterstützende- oder Pflegeassistenz sehen sie die Gefahr, dass durch diese sich der menschliche Kontakt in ihrem Alltag reduziert, weil einfache Tätigkeiten wie Füttern oder Putzen, die auch ohne Pflegekraft ausgeführt werden können, von Roboter übernommen werden [ebd.: 29]. Hier besteht die Gefahr, dass die Älteren das Gefühl einer Verobjektivierung erfahren, sowie dass sie die Kontrolle über ihr Leben zu verlieren. Die Abnahme sozialer Interaktion im Alltag führt oftmals zu einer Zunahme des Stressempfindens und zur Verringerung kognitiver Fähigkeiten, allgemein zeichnet sie sich durch eine Verschlechterung des psychischen und körperlichen Gesundheit aus [ebd.: 29/30]. Pflegebedürftigen durch das Einführen von Robotern in der Pflege zusätzlich sozialer Interaktion zu entziehen bewerten Sharkey und Sharkey als ethisch nicht vertretbar, wobei sie anmerken, dass Roboter im besten Fall durch Interaktion ihrerseits den Mangel kompensieren \cite[S. 30]{sharky}. Dieses bedarf weiterer Untersuchungen.\\Wenn die Anwendung von Robotern sich als wirtschaftlich erfolgreich herausstellt, besteht die Gefahr, dass sich ihr Einsatz trotz ethischer Unzulänglichkeiten durchsetzt. Gründe dafür können die bereits genannten demographischen Entwicklungen in Industrienationen wie auch die mangelhaften Beschäftigungszahlen im Pflegebereich sein [ebd.: 28 ff].
Problematisch ist, dass der Einsatz von Robotern im Gesundheitsbereich die Bereitstellung von  weiteren Mitteln zur Weiterentwicklung und Erforschung neuer Technologien rechtfertigt \cite[57]{haker}. Haker [ebd.] bemerkt hierzu, dass wenn die Technologien nunmal vorhanden sind, sie dementsprechend auch auf diesen neuen Bereich angewandt werden. Bevor überhaupt ethische Kriterien und Regeln für den Einsatz von Robotern diskutiert und festgelegt werden, wird ihr Einsatz bereits vorangetrieben, da der Stand der Technologie diesen hergibt. Es existiert weltweit im Bereich der Technologieforschung und -entwicklung ein großes Interesse, Roboter kontinuierlich weiterzuentwickeln und im Bereich der Anwendung zu Erproben. Zudem stellt sich die Frage, wer darüber entscheidet, wie die Roboter agieren und wie sie designed \ref{sec:values} sind. "In  Feldern  wie  der  Pflegerobotik,  die  sich  noch  in  einem  sehr  frühen  Stadium  der  Entwicklung  befinden  und  deren  Anwendungsmöglichkeiten  sich  folglich  erst  unscharf  abzeichnen,  ist  diese  […] folgenorientierte  Herangehensweise  mit  grundlegenden Problemen  konfrontiert  –  schließlich  liegen  die  Technikfolgen  noch  weitgehend  im  Dunkeln,  womit  auch  deren  ethische  Beurteilung weitgehend ins Leere läuft (was jedoch im Grunde ein Problem unzureichenden prognostischen Wissens und keines der normativen Bewertung an sich darstellt)” \cite[S. 148]{kehl}.
%\chapter{Ethische Erläuterungen zu Robotik in der Pflege}
%\label{ch:ethik}
%\section{Transparenz, Big-Data, Privacy, values in Design}
%\label{sec:trans}
%Verschiedene Veröffentlichungen beschäftigen sich in den letzten Jahren mit der Frage, inwiefern der Einsatz von Robotik in der Pflege ethisch vertretbar ist. Haker \cite[S. 56]{haker} nennt drei Bereiche, in denen neue Technologien hilfreich sein können und neuerdings erprobt werden: die Pflegeassistenz, das Monitoring sowie die soziale Förderung. Als Gründe für die Entwicklung von Robotern in der Pflege führt sie den demographischen Wandel in vielen Industrienationen, den Pflegenotstand, die Beförderung von Unabhängigkeit in der häuslichen Pflege durch Monitoring und Überwachung, sowie die Kostenreduzierung an [ebd.: 56].

%\\Auch Sharkey und Sharkey sehen als wichtige Bereiche, in denen Roboter eingesetzt werden können, die Unterstützung der Älteren, sowie der sie pflegenden; das Monitoring von Verhalten und Gesundheitszustand; sowie das Gesellschaftleisten \cite[S. 27]{sharky}. Menschen können durch Roboter emotional stimuliert werden und diese können dabei helfen “Gefühle wie Traurigkeit, Zorn und Einsamkeit zu regulieren” \cite[S. 58]{haker}. Was die ethischen Konsequenzen angeht, stimmen die meisten Veröffentlichungen hinsichtlich der sich ergebenen Problematiken überein, erkennen jedoch auch den Nutzen an. An dieser Stelle wird sich eine Analyse der sozialen Unterstützung durch Roboter beschränkt. Haker stellt drei Modelle vor, in denen sie die unterschiedlichen Argumentationslinien bezüglich der Position von Robotern in sozialen Beziehungen im Pflegebereich darlegt \cite[S. 59 ff]{haker}. Im Ausschlussmodell wird angenommen, dass zwischen Menschen und Robotern keine gleichwertige Beziehung existieren kann. Roboter werden als zu sozialer Interaktion nicht fähig gesehen, problematisch sei es, wenn Pflegebedürftige dieses durch kognitive Einschränkungen nicht mehr nachvollziehen können. “[…]So warnen Kritiker nicht nur vor der Infantilisierung der Klienten, die durch die BetreuerInnen oder Angehörigen erfolgen könnte, sondern zugleich besteht auch die Gefahr der Objektivierung, indem letztlich nämlich mit den Bedürfnissen älterer Menschen gespielt werde. Die absichtliche Simulation reziproker Interaktion durch Dritte ist von einer anderen Kategorie als die selbstbestimmte Entscheidung älterer Menschen, Roboter als soziales Spielzeug in ihr Leben einzubeziehen” [ebd.: 60]. Im Ergänzungsmodell werden Roboter vor allem für gewisse Arbeiten gedacht, die durch eine Ergänzung der Pflegeleistung diese unterstützen. Dabei wird der Einsatz von Robotern immer durch das Pflegepersonal betreut, “[…]ihr Einsatz ist klar durch den therapeutischen Rahmen gegeben, den weder die Klienten (bzw. Patienten) noch die Roboter durchbrechen können oder sollen” [ebd.: 61]. Haker wirft die Frage auf, ob in dieser Argumentationslinie das autonome Handlungsvermögen von Robotern nich unterschätzt wird \cite{haker}.\\ Im Kompensationsmodell geht es darum durch den Einsatz von Robotern die Abwesenheit anderer Interaktion zu kompensieren. Wenn soziale Interaktion mit Angehörigen oder Pflegepersonal zeitlich nicht möglich ist, sei es immer noch besser wenn Roboter den zu Pflegenden ein Wohlgefühl vermitteln und ihnen Gesellschaft leisten würden. Die Argumentationslinie des Kompensationsmodells bewertet das Steigern des Wohlbefindens als wichtiger, selbst wenn die zu Pflegenden bezüglich der Reziprozität der Beziehung getäuscht werden würden, falls sie kognitiv nicht mehr feststellen können, dass sie gerade mit einem Roboter interagieren [ebd.: 62].\\ Je nach dem, welcher Argumentationslinie gefolgt wird, wird die Annahme vertreten, dass soziale Beziehungen zu Robotern niemals gleichwertig sein können, was im schlimmsten Fall von Älteren oder Demenzkranken nicht gemerkt wird, womit sie also getäuscht werden. Oder der Einsatz von Robotern wird dennoch als positiv bewertet, da anderweitige soziale Interaktion bewusst ergänzt werden kann, oder im schlechteren Fall die Abwesenheit sozialer Interaktion durch diese andere Form der Interaktion kompensiert.\\Sparrow und Sparrow lehnen den Einsatz von Robotern nicht grundsätzlich ab, befinden ethische Bedenken jedoch als zentral. Sie weisen darauf hin, wie kompliziert die einzelnen Aufgaben im Bereich der Pflege sind, und bezweifeln, dass künstliche Intelligenz in naher Zukunft solch komplexe und verantwortungsvolle Aufgaben übernehmen werden kann \cite[S.145 ff]{sparrow}. Sparrow und Sparrow gehen davon aus, dass das zentrale menschliche Bedürfnis, geliebt zu sein und umsorgt zu werden, von Robotern nur befriedigt werden kann, wenn man der Täuschung unterliegt, dass diese gleichwertige Kreaturen mit Gefühlen sind \cite[S. 154 ff]{sparrow}. Beim Empfinden dieser positiven Emotionen ist die Beziehung zwischen zwei Lebewesen zentral. Sie ist dadurch charakterisiert, dass auch ein Gegenüber Bedürfnisse und Erwartungen hat, und dass gegenseitige Forderungen die Beziehung formen, welche sowohl unvorhersehbar, als auch manchmal anmaßend und unerwartet sein können [ebd.: 149]. “For robots to be capable even of imitating these responses successfully they would need to possess physical bodies capable of the same level of expressiveness and individuality as human bodies. Moreover, entities which do not understand the facts about human experience and mortality that make tears appropriate will be unable to fulfil this caring role” \cite[S. 154]{sparrow}. Die Täuschung, sich von einem Roboter umsorgt und geliebt zu fühlen, ist vor allem problematisch, da Illusionen wirkliche Bedürfnisse nicht dauerhaft befriedigen, wirkliche Interessen nicht dauerhaft bedienen und auch ein Wohlgefühl nicht dauerhaft illusionär, oder durch Selbsttäuschung, aufrecht erhalten werden kann [ebd.:155]. Menschen wollen nicht in illusionärer Annahme glauben, dass sie von illusionären technischen Maschinen geliebt und umsorgt werden, sondern sie haben das reale Bedürfnis von Menschen geliebt und umsorgt zu werden [ebd.]. “Insofar as robots can make people happier only when they are deceived about the robots’ real nature, robots do not offer real improvements to people’s well-being; in fact the use of robots can be properly said to harm them“ [ebd.].  Auch Haker stellt diesbezüglich fest: “Ohne diese Form der imaginativen Personalisierung des Roboters verlöre das Spiel mit ihm schnell seinen Reiz – oder umgekehrt: es wäre nicht mehr klar, was der Unterschied beispielsweise zu Stofftieren wäre, mit denen ältere Menschen genauso imaginativ soziale Interaktionen aufbauen könnten“ \cite[S. 59]{haker}. Sparrow und Sparrow vermuten, dass positive emotionale Effekte, die aus der Interaktion mit Robotern hervorgehen, sich nicht über einen langen Zeitraum halten lassen, da eine authentische Beziehung nicht aufgebaut werden kann \cite[S. 149]{sparrow}. Diese Effekte beruhen somit bloß auf der Faszination einer neuen, anregenden Erfahrung. Sharkey und Sharkey beantworten die Frage, ob ein Roboter für Ältere ein*e Gefährt*In darstellen kann, mit der Festellung, dass einsame Menschen von der Anwesenheit eines Roboters profitieren können, aber sicherlich um einiges mehr profitieren würden, wenn sie sich in einem freundlichen sozialen Umfeld befänden \cite[S. 34]{sharky}. Sie schreiben jedoch auch, dass Roboter als Gefährten nicht einfach aufgrund der stattfindenden Täuschung als umethisch gewertet werden können. Ältere könnten durchaus bewusst entscheiden, dass sie die Interaktion mit Robotern genießen, auch wenn es sich um eine illusionäre Beziehung handelt [ebd.: 36]. Auch von einer Infantilisierung kann nicht in jedem Fall gesprochen werden, es gibt durchaus Situationen wo ein älterer Mensch sich entscheidet einfach die Möglichkeit zu nutzen und das Beste daraus zu machen.
%Problematisch ist, dass der Einsatz von Robotern im Gesundheitsbereich die Bereitstellung von  weiteren Mitteln zur Weiterentwicklung und Erforschung neuer Technologien rechtfertigt %\cite[S. 57]{haker}. Haker [ebd.] bemerkt hierzu, dass wenn die Technologien nunmal vorhanden sind, sie dementsprechend auch auf diesen neuen Bereich angewandt werden. Bevor überhaupt ethische Kriterien und Regeln für den Einsatz von Robotern diskutiert und festgelegt werden, wird ihr Einsatz bereits vorangetrieben, da der Stand der Technologie diesen hergibt. Es existiert weltweit im Bereich der Technologieforschung und -entwicklung ein großes Interesse, Roboter kontinuierlich weiterzuentwickeln und im Bereich der Anwendung zu Erproben. Zudem stellt sich die Frage, wer darüber entscheidet, wie die Roboter agieren und wie sie designed sind.

%\chapter{Values and Design}

%Nissenbaum (2005) setzt sich mit der Frage auseinander, inwiefern values in design, also die durch das jeweilige Design implizierten Werte, besser in technischen und naturwissenschaftlichen Bereichen reflektiert und integriert werden können. Oftmals gibt es in diesen universitären Fakultäten wenig Überschneidungen mit den Sozial- und Geisteswissenschaften, welche eine längere Tradition der Methodologie aufweisen (ebd.: lxvi). Sie schlägt vor, dass pflichtbewusste Designer*Innen ihre values in design auf drei Wissensebenen reflektieren. Zuerst werden diese hinsichtlich des Wissensstandes in Technik und Naturwissenschaft überprüft, folgend durch eine philosophische Reflexion, und zuletzt durch eine empirische Überprüfung im Anwendungsfeld (ebd.). Durchgängig sollte dabei im Fokus stehen, dass sich alle drei Bereiche durchdringen und gegenseitig beeinflussen. Sie sollten also, während der Fokus sich in einem dieser Bereiche befindet, kontinuierlich mitgedacht werden. Während also im ersten Schritt sämtliche im jeweiligen technischen Feld vorhandenen Werte auf den Prüfstand gestellt werden um sich sinnvoll für ein Design zu entscheiden, wird im zweiten Schritt geschaut, wo Werte ihren Ursprung haben, welche Relevanzen und Bedeutungen ihnen innewohnen und wie sie normativ beeinflusst sind (ebd.: lxvii). Werte können eine moralische, eine politische und eine soziale Konnotation haben (ebd.). Werden Werte ins Design bewusst integriert, sollte im nächsten Schritt überprüft werden, ob ihre Einverleibung geglückt ist, also ob die Intentionen der*s Designer*In sich erfüllen (ebd.; lxviii).  Möglich ist, ein partizipatives Design zu erstellen, wo die Entwicklungsschritte einer Technologie oder einer Maschine jeweils rückgekoppelt sind an eine Überprüfung auf ihre Wirksamkeit im späteren Anwendungsbereich (ebd.). Nissenbaum (2005) hat mit Kolleg*Innen hierfür eine eigene Methode entwickelt, in der hierfür vier Schritte durchlaufen werden. Zunächst die Entdeckung der jeweils zugrunde liegenden Werte, folgend ihre Übersetzung in das jeweilige Feld, danach die Lösung zumeist auftretender Konflikte oder Inkontinenzen und zuletzt die Überprüfung einer erfolgreichen Einverleibung (ebd.: lxix).

%\section{Datenschutz}
%Der Einsatz von Robotern in der Pflege allgemein, aber auch in Bezug auf soziale Interaktion, birgt viele Möglichkeiten, über eine reflektierte Herangehensweise bezüglich der values in design, die ethischen Vorraussetzungen positiv zu beeinflussen. Wie Sharkey und Sharkey (2010) und Kehl (2018) feststellen, bedarf es hier einer genauen Überprüfung der Grundvorraussetzungen und eine Zusammenarbeit mit den Pflegebedürftigen. Mithilfe empirischer Untersuchungen und  partizipativer Technikgestaltung wäre es möglich, wichtige Werte von Älteren herauszufinden und aufzugreifen und somit auf technischer Seite dafür zu sorgen, dass ethische Fragen in das Design eingearbeitet werden. In der Literatur angesprochene Themen sind dabei, wie eine Kameraüberwachung in das Aussehen und Auftreten von Robotern integriert werden kann oder unter welchen Voraussetzungen gewisse private Räumen wie Bad und Schlafzimmer von Robotern betreten werden dürfen. Das Aussehen von Robotern ist eine weitere wichtige Frage, werden diese eher dem Aussehen von Lebewesen nachempfunden, oder wird ein abstraktes Äußeres bevorzugt. Je abstrakter die Erscheinung eines Roboters ausfällt, desto eher lässt sich vermuten, dass Pflegebedürftiger einer Täuschung weniger schnell unterliegen. Es wäre möglich den Roboter so zu kreieren, dass einer Vereinsamung von Pflegebedürftigen durch Warnmechanismen entgegengewirkt wird. 
%(Anknüpfungspunkt values in design) \todo{Hier den Teil zu values and design rein}Werden diese im Hinblick darauf entwickelt, um Pflegebedürftigen zu helfen, oder im Hinblick darauf, Kosten im Pflegebereich einzusparen und das Arbeitspensum des Pflegepersonals zu reduzieren \cite[S. 30]{sharky}.

%\section{Bedenken beim Einsatz von Pflegerobotern}
%Wenn die Anwendung von Robotern sich als wirtschaftlich erfolgreich herausstellt, besteht die Gefahr, dass sich ihr Einsatz trotz ethischer Unzulänglichkeiten durchsetzt. Gründe dafür können die bereits genannten demographischen Entwicklungen in Industrienationen wie auch die mangelhaften Beschäftigungszahlen im Pflegebereich sein [ebd.: 28 ff].

%Sharkey und Sharkey untersuchen den Einsatz von Robotern in der Pflege vor allem unter dem Kriterium des Wohlergehens älterer pflegebedürftiger Menschen \cite[S. 28]{sharky}. In Bezug auf die unterstützende- oder Pflegeassistenz sehen sie die Gefahr, dass durch diese sich der menschliche Kontakt in ihrem Alltag reduziert, weil einfache Tätigkeiten wie Füttern oder Putzen, die auch ohne Pflegekraft ausgeführt werden können, von Roboter übernommen werden [ebd.: 29]. Hier besteht die Gefahr, dass die Älteren das Gefühl einer Verobjektivierung erfahren, sowie dass sie die Kontrolle über ihr Leben zu verlieren. Die Abnahme sozialer Interaktion im Alltag führt oftmals zu einer Zunahme des Stressempfindens und zur Verringerung kognitiver Fähigkeiten, allgemein zeichnet sie sich durch eine Verschlechterung des psychischen und körperlichen Gesundheit aus [ebd.: 29/30]. Pflegebedürftigen durch das Einführen von Robotern in der Pflege zusätzlich sozialer Interaktion zu entziehen bewerten Sharkey und Sharkey als ethisch nicht vertretbar, wobei sie anmerken, dass Roboter im besten Fall durch Interaktion ihrerseits den Mangel kompensieren \cite[S. 30]{sharky}. Dieses bedarf weiterer Untersuchungen. Roboter in der Pflege können  positive Effekte haben, sollten sie die Abhängigkeit vom Pflegepersonal reduzieren und eine längere Selbstständigkeit im eigenen Zuhause ermöglichen [ebd.: 31].  (Siehe conclusion Sharkey2 für Regeln) \todo{sharkey anschauen}
%In ihrem Fazit bezieht Haker drei unterschiedliche ethische Rahmentheorien mit ein, um mit Hilfe dieser und der zuvor genannten Argumentationslinien Regeln für den Einsatz von Robotern in der Pflege zu erarbeiten \cite{haker}.( Regeln im Schlussteil?)\todo{schlussteil?} “Eine Ethik, die sich auf die Rahmentheorie der Menschenrechte stützt, muss klären, wie sich der Einsatz von Robotern zum Recht auf Wohlergehen, auf Privatheit, auf Gewaltfreiheit und insgesamt auf den Schutz der persönlichen Integrität verhält “ [ebd.: 63/64]. Innerhalb einer solchen Ethik stellt sich grundsätzlich die Frage, ob der Einsatz von Robotern nicht die Würde älterer Menschen verletzt. Wird im Sinne eine Befähigungsethik im Hinblick auf die Steigerung des Wohlbefinden argumentiert, stellt sich die zentrale Frage “unter welchen Bedingungen Roboter Menschen nicht nur stimulieren, sondern tatsächlich in ihren Fähigkeiten fördern und zu mehr Freiheit ermächtigen” [ebd.: 64]. Auch eine utilitaristische ethische Ausrichtung betrachtet eine Steigerung des Wohlbefindens \cite[S.8]{mill}, und somit die Effekte der Interaktion zwischen Mensch und Roboter. Diese würde allerdings argumentativ so weit gehen, dass die positiven Effekte wichtiger sind, als die menschliche Würde, sollten Menschen nicht mehr in der Lage sein eine Verletzung ihrer Würde zu reflektieren \cite[S. 64]{haker}.\\Kehl spricht sich für eine umfassende Auseinandersetzung mit der Frage, inwiefern der Einsatz von Robotik in der Pflege sinnvoll sein kann, aus \cite[S. 143 ff]{kehl}. “Für  die  Pflegerobotik  erscheint  dies  besonders  drängend, da man es hierbei mit einer Technologie zu tun hat, die aufgrund ihrer Interaktivität  und  zunehmenden  Autonomiefähigkeit  ganz  neue  technische  Möglichkeiten  bietet,  die  aber  gerade  nicht  Selbstzweck  sein  dürfen,  sondern  die  es  sinnvoll  –  und  zwar  im  Sinne  guter  Pflege  –  nutzbar  zu  machen  gilt” [ebd. 144]. Als Instrument schlägt er die Bedürfnis- und Bedarfsanalyse vor.
Sowohl was das Thema Transparenz angeht, als auch Values in Design, ist es notwendig bei der Entwicklung neuer Technologien in einen kontinuierlichen Austausch mit denjenigen zu treten, die vom Einsatz dieser Technologien letztlich betroffen sind. Dies kann einerseits die Akzeptanz fördern, andererseits das Anwendungsspektrum verbessern und die Wirkmächtigkeit verbessern.
Bei der interdisziplinären Bedarfserhebung geht es darum, “ein  genaues  Verständnis  für  die  tatsächlichen  (und  nicht  nur  unterstellten)  Lebens-  und  Problemlagen  der  potenziellen  Nutzer  zu  entwickeln” [ebd.]. Hierauf folgt eine partizipative Technikgestaltung, welche die Technikanwendung in Hinblick auf die Zielgruppe Nutzer*Innenfreundlich gestaltet.

%"In  Feldern  wie  der  Pflegerobotik,  die  sich  noch  in  einem  sehr  frühen  Stadium  der  Entwicklung  befinden  und  deren  Anwendungsmöglichkeiten  sich  folglich  erst  unscharf  abzeichnen,  ist  diese  […] folgenorientierte  Herangehensweise  mit  grundlegenden Problemen  konfrontiert  –  schließlich  liegen  die  Technikfolgen  noch  weitgehend  im  Dunkeln,  womit  auch  deren  ethische  Beurteilung weitgehend ins Leere läuft (was jedoch im Grunde ein Problem unzureichenden prognostischen Wissens und keines der normativen Bewertung an sich darstellt)”. (Kehl 2018: 148) \cite[S. 148]{kehl} \\\\\\\



Allerdings  hat  die  Ethik  aus  
148C. Kehlverschiedenen  Gründen  Mühe,  in  Problemkonstellationen,  wie  sie  sich  aktuell  im  Bereich der Pflegerobotik abzeichnen, dieser Aufgabe gerecht zu werden:1.Unsicherheit  der  Folgenabschätzung:  Quelle  normativer  Unsicherheiten  und  mithin  Gegenstand der ethischen Bewertung ist in der Regel nicht die Technik an sich, sondern deren  Rolle  in  konkreten  soziotechnischen  Handlungszusammenhängen  (Grunwald  2013a, S.4). Dazu gehören wesentlich die intendierten sowie nicht intendierten Folgen, die  mit  einem  Technikeinsatz  verbunden  sind  und  deren  moralische  Implikationen  es  ethisch  zu  klären  gilt.  In  Feldern  wie  der  Pflegerobotik,  die  sich  noch  in  einem  sehr  frühen  Stadium  der  Entwicklung  befinden  und  deren  Anwendungsmöglichkeiten  sich  folglich  erst  unscharf  abzeichnen,  ist  diese  konsequenzialistische  (folgenorientierte)  Herangehensweise  mit  grundlegenden 



VALUES IN DESIGN!!!

\chapter{Fazit}
\label{ch:fazit}

In ihrem Fazit bezieht Haker drei unterschiedliche ethische Rahmentheorien mit ein, um mit Hilfe dieser und der zuvor genannten Argumentationslinien Regeln für den Einsatz von Robotern in der Pflege zu erarbeiten \cite{haker}. “Eine Ethik, die sich auf die Rahmentheorie der Menschenrechte stützt, muss klären, wie sich der Einsatz von Robotern zum Recht auf Wohlergehen, auf Privatheit, auf Gewaltfreiheit und insgesamt auf den Schutz der persönlichen Integrität verhält “ [ebd.: 63/64]. Innerhalb einer solchen Ethik stellt sich grundsätzlich die Frage, ob der Einsatz von Robotern nicht die Würde älterer Menschen verletzt. Wird im Sinne eine Befähigungsethik im Hinblick auf die Steigerung des Wohlbefinden argumentiert, stellt sich die zentrale Frage “unter welchen Bedingungen Roboter Menschen nicht nur stimulieren, sondern tatsächlich in ihren Fähigkeiten fördern und zu mehr Freiheit ermächtigen” [ebd.: 64]. Auch eine utilitaristische ethische Ausrichtung betrachtet eine Steigerung des Wohlbefindens \cite[8]{mill}, und somit die Effekte der Interaktion zwischen Mensch und Roboter. Diese würde allerdings argumentativ so weit gehen, dass die positiven Effekte wichtiger sind, als die menschliche Würde, sollten Menschen nicht mehr in der Lage sein eine Verletzung ihrer Würde zu reflektieren \cite[64]{haker}. Kehl spricht sich für eine umfassende Auseinandersetzung mit der Frage, inwiefern der Einsatz von Robotik in der Pflege sinnvoll sein kann, aus \cite[143 ff]{kehl}. \“Für  die  Pflegerobotik  erscheint  dies  besonders  drängend, da man es hierbei mit einer Technologie zu tun hat, die aufgrund ihrer Interaktivität  und  zunehmenden  Autonomiefähigkeit  ganz  neue  technische  Möglichkeiten  bietet,  die  aber  gerade  nicht  Selbstzweck  sein  dürfen,  sondern  die  es  sinnvoll  –  und  zwar  im  Sinne  guter  Pflege  –  nutzbar  zu  machen  gilt” [ebd. 144]. Als Instrument schlägt er die Bedürfnis- und Bedarfsanalyse vor. Sowohl was das Thema Privatsphäre angeht, als auch values in design, ist es notwendig bei der Entwicklung neuer Technologien in einen kontinuierlichen Austausch mit denjenigen zu treten, die vom Einsatz dieser Technologien letztlich betroffen sind [ebd.:].% Dies kann einerseits die Akzeptanz fördern, andererseits das Anwendungsspektrum verbessern und die Wirkmächtigkeit verbessern. Bei der interdisziplinären Bedarfserhebung geht es darum, “ein  genaues  Verständnis  für  die  tatsächlichen Lebens-  und  Problemlagen  der  potenziellen  Nutzer  zu  entwickeln” [ebd.]. Hierauf folgt eine partizipative Technikgestaltung, welche die Technikanwendung in Hinblick auf die Zielgruppe Nutzer*Innenfreundlich gestaltet.
Häufig  wird  der  Kritikpunkt  vorgebracht,  dass  Roboter  keine  echte  Zuwendung  ersetzen  können.  Es wurde festgestellt,  dass  Roboter  zwar  zur  Stimmungs-aufhellung   von   Demenzpatienten   beitragen   können,   aber   nicht   die   Wirkung   von   zwischenmenschlichen  Kontakten  erzielen \cite[79]{vier}. Es könnte eine Abnahme sozialer Interaktion stattfinden, Pflegebedürftige könnten zunehmend vernachlässigt werden, wobei die Vernachlässigung möglicherweise damit entschuldigt wird, dass Roboter emotionale und körperliche Bedürfnisse bedienen. Es gibt aber auch Evidenz  für  den  therapeutischen  Nutzen  von  simulierten  sozialen  Stimuli. Demenzkranke Patienten könnten so vom Roboter angeregt werden, mit ihrem sozialen Umfeld zu interagieren \cite[79]{vier}. Wenn ein Roboter primär eingesetzt wird, um Gesellschaft zu leisten, wäre es möglich, dass er registriert, wieviel Zeit Pflegebedürftige ohne soziale Interaktion verbracht haben. Die Roboter könnten eine Warnung abgeben, wenn die Pflegebedürftigen zu lange Zeiträume allein verbracht haben. Roboter in der Pflege können  positive Effekte haben, sollten sie die Abhängigkeit vom Pflegepersonal reduzieren und eine längere Selbstständigkeit im eigenen Zuhause ermöglichen \cite[31]{sharky}. Dem zufolge kann der Einsatz von Robotern in der Pflege zwar zu einem Zeitgewinn führen, es  besteht  jedoch  die  akute  Gefahr,  dass  die  gewonnenen  Minuten  aus  ökonomischen  Gründen nicht  dem  einzelnen  Zentrumsbewohner  zugutekommen.  Vielmehr  dürfte  der  Robotereinsatz  dazu  führen,  dass  eine  höhere  Anzahl  an  Patienten  von  der  gleichen  Anzahl an Pflegekräften oder dieselbe Anzahl an Patienten von einer geringeren Anzahl an  Pflegekräften  betreut  wird\cite[143]{sparrow}. Die in Sektion \ref{sec:priv} dargelegten Erläuterungen könnten einen Verlust von Privatsphäre und Einschränkung der persönlichen Freiheit bedeuten.
Täuschung und Infantilisierung könnten vorkommen, wenn Pflegebedürftige dazu angeregt werden mit Robotern eine gleichwertige Beziehung aufzubauen Wenn die Kontrolle über in der Hand der Pflegebedürftigen liegt, sind, sollte etwas schief gehen, die Verantwortlichkeiten unklar.





%Im Gegensatz zu Computer Ethics allgemein, wirft der Einsatz von Robotern weitergehende Fragen auf. Roboter sind mobil und können Individuen verfolgen. Sie verkörpern in ihrem Äußeren Lebewesen. Ihre Erscheinungsform kann dazu führen, dass sie auch in Bereichen akzeptiert werden, in denen Videoüberwachung unter anderen Bedingungen nicht vorstellbar wäre (siehe \ref{sec:values}. Ihre den Lebewesen nachempfundene Erscheinungsform kann zudem bewirken, dass den Robotern soziale Fähigkeiten zugetraut werden, welche sie nicht besitzen.


Ein Bereich, der durch Richtlinien von Sharkey und Sharkey geregelt werden könnte, ist die Zeitdauer, die Pflegebedürftige allein, nur in Begleitung eines Roboters zurückgelassen werden dürfen \cite[38]{sharky}. Es könnte eine Vorschrift geben, dass jeder Roboter erst fragen muss, ob körperlicher Kontakt akzeptiert wird, bevor ein Roboter einen Menschen stützt oder anhebt.
Weitere Richtlinien sollten durch umfassende Auswertungen der Auswirkungen eines Robotereinsatzes auf Pflegebedürftige erstellt werden.











   


REGELN!!!!!

%Regeln (Haker 2014: 65/66)
%• Jedem Menschen steht es frei, sich einen „sozialen Roboter“ zuzulegen – dies gilt auch für Menschen im Alter. Sofern sie diese kontrollieren können, spricht nichts gegen die spielerische, imaginative Interaktion von Menschen und (sozialen) Robotern. 
%• Falls in Betreuungskonstellationen Interaktionen zwischen Klienten/Patienten und BetreuerInnen bzw. Mensch-Tier-Interaktionen zeitlich und räumlich möglich sind, diese aber Gefahr laufen, durch den Einsatz von Robotern reduziert zu werden, ist den spontanen reziproken Interaktionen der Vorzug zu geben. Eine Ergänzung bzw. das Spiel mit sozialen Robotern muss daher notwendig in die allgemeinen (menschlichen) Interaktionen integriert werden. 
Ein infantilisierender Einsatz von Robotern ist auszuschließen. Sofern ältere Menschen dazu in der Lage sind, müssen sie über die Funktionsweise von Robotern informiert werden. Sofern sie dazu nicht in der Lage sind, müssen Dritte darin geschult werden, Roboter ergänzend und nicht kompensativ einzusetzen. Dies erfordert eine situations- und kontextabhängige Schulung. 

Sofern Roboter Handlungen ausführen, die bestimmte Reaktionen hervorrufen, ist die Intimität und Integrität der Personen zu gewährleisten. Bewusste Täuschungen sind auszuschließen, was durch Informationen und Kommunikation sowie über die Integration der Roboter in die sozialtherapeutische Begleitung zu gewährleisten ist. 

Roboter können unter Umständen ein Sicherheitsrisiko darstellen. Die Verantwortung für ihren Einsatz ist zu klären, insbesondere dann, wenn sie in Abwesenheit von Betreuungspersonen eingesetzt werden. 



%\include{kapitel5}
%\include{kapitel6}
%\include{kapitel7}
%\cleardoublepage

% VERZEICHNISSE (Abbildungen, Tabellen)
% Literatur 

\nocite{Poole/Mackworth}
\nocite{cormen}
\nocite{at1}
\nocite{at2}
\nocite{at3}
\nocite{at4}
\bibliographystyle{plainnat}
\bibliography{Bachelorarbeit}
\cleardoublepage

% ERKLÄRUNG
%\chapter*{Eidesstattliche Versicherung}
\thispagestyle{empty}
\addcontentsline{toc}{chapter}{Eidesstattliche Versicherung}

Hiermit versichere ich an Eides statt, dass ich die vorliegende Arbeit selbstständig und ohne fremde Hilfe angefertigt und mich anderer als der im beigefügten Verzeichnis angegebenen Hilfsmittel nicht bedient habe.
Alle Stellen, die wörtlich oder sinngemäß aus Veröffentlichungen entnommen wurden, sind als solche kenntlich gemacht. 


\vspace{2cm} 

\noindent Hamburg, den \uline{~~~~~~~~~~~~~~~~~~~~}~~~~~Unterschrift: \uline{~~~~~~~~~~~~~~~~~~~~~~~~~~~~~~~~~~~~~~~~~~~~~~~~~~} 

    
\end{document}